% Preamble — pacotes e configurações comuns.
% Baseado em 00-pacotes.tex do modelo ufpr-abntex (https://github.com/eekBR/ufpr-abntex),
% acrescido dos pacotes usados neste trabalho. Carregado depois de \usepackage{UFPR}.

% ----------------------------------------------------------
% Pacotes básicos (modelo UFPR)
% ----------------------------------------------------------
\usepackage[T1]{fontenc}
\usepackage[utf8]{inputenc}
\usepackage{lastpage}
\usepackage{indentfirst}
\usepackage{microtype}
\usepackage{ifthen}
\usepackage[brazil]{babel}
\usepackage[final]{pdfpages}    % termo de aprovação e ficha catalográfica em PDF
\usepackage{amsmath, amsfonts, amssymb}
\usepackage{ragged2e}

% ----------------------------------------------------------
% Pacotes do trabalho
% ----------------------------------------------------------
\usepackage{graphicx}
\usepackage{booktabs}
\usepackage{siunitx}
\usepackage{xcolor}
\usepackage{csquotes}
\usetikzlibrary{shapes,arrows.meta,positioning,calc}

% ----------------------------------------------------------
% Citações e referências — ABNT (biblatex-abnt), como no modelo UFPR
% ----------------------------------------------------------
\usepackage[
    style=abnt,
    backend=biber,
    repeatfields=true,
    url=true,
    maxbibnames=99,
    mincitenames=1,
    maxcitenames=2,
    hyperref=true,
    giveninits=true,
    uniquename=false,
    uniquelist=false,
    natbib=true,
]{biblatex}

\addbibresource{refs.bib}

% NBR 10520:2023 — sobrenomes em minúsculas nas chamadas
\renewcommand*{\mkbibnamefamily}[1]{#1}

% Espaçamento simples dentro de cada referência, alinhamento à esquerda
\renewcommand*{\bibsetup}{%
  \renewcommand{\baselinestretch}{1.0}%
  \RaggedRight
  \interlinepenalty=5000\relax
  \setlength{\parskip}{0pt}%
  \setlength{\itemsep}{0pt}%
  \setlength{\parsep}{0pt}%
  \setlength{\bibitemsep}{0.5\baselineskip}%
  \setlength{\bibparsep}{\parskip}%
}

\DefineBibliographyStrings{brazil}{%
  urlfrom = {Dispon\'ivel em},
}

% Título REFERÊNCIAS sem numeração, presente no sumário
\defbibheading{bay}[\bibname]{%
  \chapter*{#1}%
  \markboth{#1}{#1}%
  \addcontentsline{toc}{chapter}{\bibname}%
}

% URL sem < > (NBR 6023:2018)
\DeclareFieldFormat{url}{\bibstring{urlfrom}\addcolon\space\url{#1}}
\DeclareFieldFormat{urldate}{\bibstring{urlseen}\addcolon\addspace #1}

% ----------------------------------------------------------
% Sumário e listas: espaçamento e caixa alta (modelo UFPR)
% ----------------------------------------------------------
\makeatletter
\pretocmd{\chapter}{\addtocontents{toc}{\protect\addvspace{-12\p@}}}{}{}
\pretocmd{\section}{\addtocontents{toc}{\protect\addvspace{-3\p@}}}{}{}
\makeatother

% Seções secundárias em caixa alta no sumário; títulos de figuras e tabelas
% em caixa alta nas respectivas listas.
\let\oldcontentsline\contentsline
\def\contentsline#1#2{%
  \ifboolexpr{ test {\ifstrequal{#1}{section}}
            or test {\ifstrequal{#1}{figure}}
            or test {\ifstrequal{#1}{table}} }
    {\oldcontentsline{#1}{\MakeTextUppercase{#2}}}%
    {\oldcontentsline{#1}{#2}}%
}

% ----------------------------------------------------------
% Numeração de páginas: canto superior direito, fonte menor (modelo UFPR)
% ----------------------------------------------------------
\makepagestyle{chapfirst}
\makeoddhead{chapfirst}{}{}{\footnotesize\thepage}

\makepagestyle{simplestextual}
\makeevenhead{simplestextual}{}{}{\footnotesize\thepage}
\makeoddhead{simplestextual}{}{}{\footnotesize\thepage}
\makeevenfoot{simplestextual}{}{}{}
\makeoddfoot{simplestextual}{}{}{}

% ----------------------------------------------------------
% Ilustrações e tabelas (normas UFPR): título acima, em caixa alta, fonte 10,
% no formato ``FIGURA 1 -- TÍTULO''; linha FONTE abaixo.
% ----------------------------------------------------------
\usepackage{caption}
\DeclareCaptionTextFormat{caixaalta}{\MakeTextUppercase{#1}}
\DeclareCaptionLabelSeparator{travessao}{ -- }
\captionsetup{
    font=footnotesize,
    labelsep=travessao,
    textformat=caixaalta,
    justification=centering,
    position=top,
    skip=6pt,
}

% Linha FONTE inserida automaticamente ao final de toda figura e tabela.
% Dentro de um float específico:
%   \renewcommand{\textofonte}{Adaptado de \textcite{chave}}  — outra fonte
%   \semfonte                                                — suprime a linha
%   \notaufpr{Texto da nota}                                    — linha NOTA após a FONTE
\newcommand{\textofonte}{O autor (\imprimirdata)}
\newcommand{\textonota}{}
\newcommand{\notaufpr}[1]{\renewcommand{\textonota}{#1}}
\newcommand{\linhafonte}{\par\vspace{2mm}{\footnotesize FONTE: \textofonte.\par
  \ifdefempty{\textonota}{}{NOTA: \textonota\par}}}
\newcommand{\semfonte}{\let\linhafonte\relax}
\AtEndEnvironment{figure}{\linhafonte}
\AtEndEnvironment{table}{\linhafonte}

% ----------------------------------------------------------
% Lista de siglas (pacote glossaries; carregar após hyperref)
% ----------------------------------------------------------
\usepackage[acronym, nomain, nopostdot, nonumberlist]{glossaries}
\makenoidxglossaries
\setacronymstyle{long-short}
\renewcommand*{\glossarysection}[2][]{\pretextualchapter{#2}}
% Definição das siglas (\newacronym). 1º arg = chave; 2º = forma curta; 3º = forma longa.
% O campo description é o que aparece na Lista de Siglas; para siglas conceituais,
% inclui a referência seminal via \parencite.

% --- Conceitos e métodos ---
\newacronym{ia}{IA}{inteligência artificial}
\newacronym[description={\textit{large language model} (grande modelo de linguagem)}]%
  {llm}{LLM}{\textit{large language model}}
\newacronym[description={\textit{machine learning} (aprendizado de máquina)}]%
  {ml}{ML}{\textit{machine learning}}
\newacronym[description={\textit{skill-biased technical change} --- mudança técnica enviesada por qualificação \parencite{tinbergen1974, katzMurphy1992}}]%
  {sbtc}{SBTC}{\textit{skill-biased technical change}}
\newacronym[description={\textit{routine-biased technological change} --- mudança técnica enviesada por rotina \parencite{autorLevyMurnane2003, autorDorn2013}}]%
  {rbtc}{RBTC}{\textit{routine-biased technological change}}
\newacronym{mqo}{MQO}{mínimos quadrados ordinários}
\newacronym[description={\textit{regression discontinuity design} (regressão descontínua)}]%
  {rdd}{RDD}{\textit{regression discontinuity design}}
\newacronym{ic}{IC}{intervalo de confiança}

% --- Técnicas, formatos e protocolos ---
\newacronym[description={\textit{Application Programming Interface} (interface de programação de aplicações)}]%
  {api}{API}{\textit{Application Programming Interface}}
\newacronym[description={\textit{Digital Object Identifier} (identificador de objeto digital)}]%
  {doi}{DOI}{\textit{Digital Object Identifier}}
\newacronym{prisma}{PRISMA}{\textit{Preferred Reporting Items for Systematic Reviews and Meta-Analyses}}
\newacronym{cbo}{CBO}{Classificação Brasileira de Ocupações}

% --- Instituições ---
\newacronym{nber}{NBER}{National Bureau of Economic Research}
\newacronym[description={\textit{Institute of Labor Economics} (\textit{Institut zur Zukunft der Arbeit})}]%
  {iza}{IZA}{Institute of Labor Economics}
\newacronym{cepr}{CEPR}{Centre for Economic Policy Research}
\newacronym[description={\textit{Bank for International Settlements} (Banco de Compensações Internacionais)}]%
  {bis}{BIS}{Bank for International Settlements}
\newacronym[description={\textit{Organisation for Economic Co-operation and Development} (OCDE)}]%
  {oecd}{OECD}{Organisation for Economic Co-operation and Development}
\newacronym{ipea}{IPEA}{Instituto de Pesquisa Econômica Aplicada}
\newacronym{bcb}{BCB}{Banco Central do Brasil}
\newacronym{fgv}{FGV}{Fundação Getulio Vargas}


% ----------------------------------------------------------
% Comandos auxiliares
% ----------------------------------------------------------
\newcommand{\janelaUm}{2013--2017}
\newcommand{\janelaDois}{2018--2022}
\newcommand{\janelaTres}{2023--2026}
